\section{Introduction}\label{sec:introduction}

The closed-geodesic infinitude problem asks whether every closed Riemannian
manifold of dimension at least two has infinitely many geometrically distinct
closed geodesics. A closed geodesic is \emph{prime} if it traverses its
periodic trajectory once. Geometric distinction means distinction of images:
changing the initial point, reversing orientation, or taking an iterate
produces no new image. A prime geodesic may have self-intersections.
We resolve the sphere case of the infinitude problem.

\begin{theorem}\label{thm:spheres}
For every integer $n\ge2$ and every smooth Riemannian metric on the standard
smooth sphere $S^n$, there are infinitely many prime closed geodesics with
pairwise different images.
\end{theorem}

Birkhoff developed a minimax construction for closed geodesics on surfaces
of genus zero~\cite[Section~17]{Birkhoff1917}.
The existence of a nonconstant closed geodesic on every positive-dimensional
closed Riemannian manifold goes back to Lyusternik and
Fet~\cite{LyusternikFet1951}.
Infinitude requires separating new geodesics from repeated traversals of
old ones. Bott's iteration theory~\cite{Bott1956Iteration} and the local
Morse theory of Gromoll and Meyer~\cite{GromollMeyer1969} make that
distinction quantitative. In particular, for a simply connected closed
manifold, unbounded Betti numbers of the free loop space over a field
force infinitely many geometrically distinct closed geodesics.
Vigu\'e-Poirrier and Sullivan~\cite{ViguePoirrierSullivan1976} showed that
rational cohomology requiring at least two algebra generators supplies
this topological growth. Spheres have one-generator cohomology and bounded
degreewise loop homology, so this criterion does not settle their infinitude
problem.

The two-sphere case is classical. It follows from Bangert's work together
with that of Franks or Hingston~\cite{Bangert1993,Franks1992,Hingston1993}.
De Philippis, Marini,
Mazzucchelli and Suhr extended this result to reversible Finsler
two-spheres~\cite[Theorem 1.5]{DePhilippisMariniMazzucchelliSuhr2022}.
Albach's more recent work gives a growth exponent of at least two for
the number of geodesic images of bounded length in that reversible
Finsler setting~\cite[Theorem 1.1]{Albach2026}.

In higher dimensions, Rademacher proved infinitude for strongly bumpy
metrics on compact simply connected manifolds and proved that such metrics
are generic~\cite[Theorems 1.1--1.3]{Rademacher1994}.
Finite multiplicity results also hold without genericity: Long and Duan
proved the existence of at least two geometrically distinct closed
geodesics on every Riemannian three-sphere~\cite{LongDuan2009}.
These results do not establish infinitude for arbitrary degenerate metrics;
a limiting argument from generic metrics need not preserve distinct prime images.
For manifolds with positive first Betti number, Contreras and Mazzucchelli
recently proved infinitude for generic metrics
\cite[Corollary~B]{ContrerasMazzucchelli2025}.
Our higher-dimensional argument addresses the unrestricted metric problem.

Earlier infinitude claims also require care. Asselle and
Mazzucchelli~\cite[Introduction]{AsselleMazzucchelli2018} identify the proof in
Klingenberg's monograph~\cite{Klingenberg1978} as erroneous.
Charles~\cite[Theorem 6.14]{Charles2019} later claimed infinitude on every
closed Riemannian manifold. We identify two unsupported steps in that argument:
the boundary-rank equality in Equation~(53) requires maximal ranks not
established there, and the constant-ball modification in Equation~(54)
does not provide the nondegenerate Morse critical points used in the
conclusion. These are objections to the proofs, not counterexamples to
their conclusions.

\begin{corollary}\label{cor:finite-cover}
Let $N$ be a closed manifold admitting a finite smooth cover diffeomorphic to
a standard sphere $S^n$, where $n\ge2$. Every smooth Riemannian metric on $N$
has infinitely many geometrically distinct prime closed geodesics.
\end{corollary}

This includes real projective spaces and smooth spherical space forms.
The passage between a manifold and its finite covers is classical
\cite[pp.~633--634]{ViguePoirrierSullivan1976}.
Section~\ref{sec:covers} proves the stated form by showing that projection
has finite fibers on prime geodesic images, including self-intersections
and nonregular covers. The conclusion concerns the number
of images. Even on a round sphere, all prime geodesics have the same length.

The finite-cover corollary also combines with a theorem of Rademacher and
Taimanov to settle the infinitude problem in dimension three.

\begin{corollary}[Closed three-manifolds]\label{cor:three-manifolds}
Let $M$ be a nonempty closed smooth three-manifold, not necessarily orientable.
Every smooth Riemannian metric on $M$ has infinitely many prime closed
geodesics with pairwise distinct images.
\end{corollary}

\begin{proof}
It suffices to work on one connected component, so assume that $M$ is
connected. If $\pi_1(M)$ is infinite, Rademacher and Taimanov
\cite[Theorem~2]{RademacherTaimanov2018} give infinitely many geometrically
distinct closed geodesics for the given metric. Their Riemannian convention
identifies geodesics exactly when their images coincide
\cite[Section~3]{RademacherTaimanov2018}. Taking the least positive period
gives a prime representative with the same image. Their theorem includes
the nonorientable case.

If $\pi_1(M)$ is finite, the elliptization case of Perelman's geometrization
theorem says that $M$ is diffeomorphic to a spherical space form
$S^3/\Gamma$, where $\Gamma$ is finite and acts freely and smoothly on $S^3$
\cite[Corollary~0.2(b)]{MorganTian2007}. Composing the quotient map with a
diffeomorphism $S^3/\Gamma\to M$ gives a finite smooth cover $S^3\to M$.
Corollary~\ref{cor:finite-cover} therefore applies. Only the smooth covering
is used, so the given Riemannian metric remains arbitrary.
\end{proof}

This corollary uses only the infinitude conclusion from Rademacher--Taimanov.
It carries no quantitative length-counting assertion for the
finite-fundamental-group case.

\subsection{Idea of the proof}

Assume that a metric on $S^n$, $n\ge3$, has only finitely many prime geodesic
images. Every nonconstant critical loop is then an iterate of one of finitely many
oriented trajectories. Their local homology has uniformly bounded rank and
lies near finitely many lines relating index to length. The difficulty is
that sublevel homology can also contain classes that disappear at a later
level. We first control these temporary classes and then use the symmetry
of double loops to obtain a contradiction.

The ordinary homological input is a uniform filling statement for every
metric, independent of the finite-prime assumption. Square-root energy
is the $L^2$-norm of the velocity of a period-one loop.
There is a constant $D$ such that a cycle below square-root energy $a$
which bounds in the full loop space already bounds below $a+D$,
in every degree (Theorem~\ref{thm:ordinary-death}). Fixed-degree shortening
first gives a bound that may depend on degree. A cut operation paired with
a cohomology power lowers a high degree into a fixed finite range, and a
loop product restores it. The resonance theorem of Hingston and
Rademacher~\cite{HingstonRademacher2013} makes the lengths removed and added
differ by a bounded amount.

The operations underlying this step come from Chas and Sullivan's loop
product~\cite{ChasSullivan1999} and Goresky and Hingston's relative
cohomology product and its length filtration~\cite{GoreskyHingston2009}.
Cohen and Jones gave a Thom-spectrum realization of the loop product
and its unit~\cite{CohenJones2002}.
Quantitative shortening of loop families has a related history in the
work of Nabutovsky and Rotman~\cite{NabutovskyRotman2013}; the
fixed-degree additive filling estimate of Block, Manin and Weinberger
\cite[Theorem 5.1]{BlockManinWeinberger2025} is a direct predecessor of
the first step here. We prove that step in the strong $H^1$ topology
before removing its degree dependence. Recent work of Cieliebak, Oancea
and Shelukhin extends resonance to further compact rank-one symmetric
spaces under dimension and coefficient restrictions, and to finite sphere
quotients with local coefficients
\cite[Theorem~6.1 and Corollary~8.4]{CieliebakOanceaShelukhin2026}.

The essential comparison is the filtered cut--join identity of
Theorem~\ref{thm:filtered-interchange}. Its chain constructions build on
the string-topology models of Hingston and Wahl~\cite{HingstonWahl2023}.
The comparison uses an absolute cut with a diagonal-relative Thom cocycle;
its value on a suitable power is the constant-loop fundamental class,
the unit for the loop product (Proposition~\ref{prop:absolute-unit}).
We prove both the identity and this unit statement, including their
filtration bounds.

The remaining argument uses reflection of loops and translation by half a
period. With coefficients in $\mathbb F_2$, let $B_a$ denote
reflection-equivariant cohomology below
square-root energy $a$, relative to constant loops, and let $B_\infty$
denote the corresponding global group. Define
\[
 T(a)=\dim\coker(B_\infty\longrightarrow B_a),
\]
where the dimension is summed over all degrees. Thus $T(a)$ counts the
extra reflection classes at that sublevel. Theorem~\ref{thm:bounded-transients}
shows that $T(a)$ is uniformly bounded under the finite-prime assumption.
The proof compares two filtrations on the half-translation fixed locus:
the actual equivariant degree and the degree of its half-length root.
An exact-sequence calculation bounds their total discrepancy by $O(a)$
(Theorem~\ref{thm:finite-cut-budget}). A long string of extra classes would
give $\Omega(a^2)$ discrepancy. The ordinary filling estimate excludes the
remaining possibilities.

Finally, the full orthogonal-group action gives
\[
                         T(2a)\ge T(a).
\]
At each stage of the round-metric filtration, the global basis has a
lower-degree family, called the bottoms, and a second family obtained by
cup product with the sphere class pulled back by evaluation at time zero.
The equality case requires every nonzero bottom restriction from an odd
round-filtration stage to have a partner whose restriction is independent
modulo all bottom restrictions in its degree
(Theorem~\ref{thm:doubling}). Because $T$ is a bounded
integer, an interval where it attains its maximum propagates to arbitrarily
large intervals under doubling. A bottom generator must first appear at
some crossing in such an interval. Evaluation cup vanishes on that local
crossing, so its partner is still zero. This contradicts the equality case.

\subsection{Organization and reusable ingredients}

Section~\ref{sec:setup} fixes the loop-space, coefficient and symmetry
conventions. Section~\ref{sec:local} establishes the local crossing theorem with one
common flow defining all restriction maps. Its rank proof reduces arbitrary
degeneracy to finitely many root types; it also includes primes of zero mean
index. Sections~\ref{sec:shortening}--\ref{sec:ordinary} prove uniform ordinary
filling. Section~\ref{sec:equivariant} computes the global generators and
their lifts by natural Gysin sequences and the round filtration.
Sections~\ref{sec:budget}--\ref{sec:activation} give the finite-cut estimate,
the transient bound and the contradiction.

Two features are useful beyond the final rank argument. The filtered
cut--join comparison turns fixed-degree filling into a bound independent
of degree. The finite-cut comparison controls
actual restriction images before localization has become an isomorphism.
These statements are isolated with their assumptions and outputs so that
their uses do not depend on the surrounding contradiction argument.
